% 論文節の日本語訳（読み物 PDF）用 preamble
% lualatex + jlreq
\documentclass[12pt,paper=a4]{jlreq}
\usepackage{luatexja-fontspec}
\usepackage{amsmath,amssymb,bm}
\usepackage{graphicx}
\usepackage{hyperref}
\usepackage{enumitem}
\usepackage{booktabs}
\usepackage{microtype}
\usepackage{xcolor}
\usepackage{caption}
\usepackage{tcolorbox}
\tcbuselibrary{breakable,skins}

\hypersetup{
  colorlinks=true,
  linkcolor=blue!60!black,
  urlcolor=blue!60!black,
  citecolor=blue!60!black,
}

\graphicspath{{images/}{./}}

\setlength{\parindent}{1em}
\setlength{\parskip}{0.55em}
\setlist[itemize]{leftmargin=1.6em,itemsep=0.25em,topsep=0.35em}
\setlist[enumerate]{leftmargin=1.8em,itemsep=0.25em,topsep=0.35em}

\captionsetup{
  font=small,
  labelfont=bf,
  skip=0.4em,
  justification=raggedright,
  singlelinecheck=false,
}

\newcommand{\vx}{\bm{x}}
\newcommand{\vu}{\bm{u}}
\newcommand{\veps}{\bm{\epsilon}}
\newcommand{\vw}{\bm{w}}

\DeclareMathOperator*{\argmax}{arg\,max}
\DeclareMathOperator*{\argmin}{arg\,min}
\DeclareMathOperator{\Tr}{Tr}
\DeclareMathOperator{\diag}{diag}

% 命題・定理ボックス（ページ跨ぎ時も切れ目を抑える）
\newtcolorbox{propositionbox}[1][]{%
  enhanced,
  breakable,
  colback=blue!4!white,
  colframe=blue!45!black,
  coltitle=white,
  fonttitle=\bfseries\small,
  title={#1},
  attach boxed title to top left={yshift=-2mm, xshift=3mm},
  boxed title style={
    colback=blue!45!black,
    colframe=blue!45!black,
    boxrule=0pt,
    arc=1pt,
  },
  boxrule=0.5pt,
  arc=1.5pt,
  left=5pt, right=5pt, top=8pt, bottom=5pt,
  before skip=0.9em,
  after skip=0.9em,
  pad at break=3mm,
}

% 図はフロートせず本文位置に置く（分断を防ぐ）
\newcommand{\translationfigure}[2]{%
  \begin{center}
    \includegraphics[width=0.88\linewidth]{#1}\\[0.35em]
    {\small #2\par}
  \end{center}
  \vspace{0.4em}%
}
